For a more complex model, we consider an application of
``multilevel regression and post--stratification,'' or MrP 
\citep{park:2004:mrp}.  We studied an example from \citet{lopezmartin:2022:mrptutorial}
in which a subsample of a national survey is used to infer a prediction
on the entire population.  Since the textbook example involves simulated subsampling,
a ground truth is available for comparison.

MrP models can be motivated by the following survey sampling problem. Let $\zn$
denote a vector of regressors, and $\yn$ a binary survey response.  Suppose we
observe $N_{\survey}$ observations of the form $\x_n = (\z_n, \y_n)$ from a
``survey population,'' and $N_{\target}$ observations consisting only of $\z_m$
from a ``target population.''  We assume that $\expect{\p(\yn \vert \zn)}{\yn}$
is the same in both the survey and target populations, but that the marginal
distributions of $\zn$ may differ, perhaps due to biased survey sampling.  Let
$\p_\target(\cdot)$ denote the target population distribution.  We would like to
use the observed data to estimate $\expect{\p_\target(\yn)}{\yn} =
\expect{\p_\target(\zn)}{\expect{\p(\yn \vert \zn)}{\yn}}$ in the target
population.  The approach of MrP is to use the survey data to form a regression
estimate of $\zn \mapsto \expect{\p(\yn \vert \zn)}{\yn}$ (the ``multilevel
regression'' part of MrP) and then to average the regression's predictions using
the target population regressors (the ``post--stratification'' part of MrP).

A typical MrP model assumes that the expectation of the response conditional on
the regressors is given a finitely parameterized regression model.  For example,
one might specify a logistic regression model for a binary survey response.  The
regression model gives an explicit expression for $\expect{}{\y | \z, \theta}$
which, by assumption, holds in both the survey and target populations.
Typically, the MrP model also specifies hierarchical, or ``multilevel'' prior
structure on the parameters, possibly including unknown hyperparameters. For
compactness of notation we assume that these unknown hyperparameters are also
included in $\theta$, even if they are not used directly in the computation of
$\expect{}{\y | \z, \theta}$.  For any particular value of $\theta$ we can form
the ``poststratified'' estimator of $\expect{\p_\target(\yn)}{\yn}$, which is
our quantity of interest:
%
\begin{align}
    \g(\theta) = \frac{1}{N_{\target}} \sum_{m=1}^{N_{\target}} \expect{}{\yn | \z_m, \theta}.
    \eqlabel{mrpg}
\end{align}
%
The expression for $\g(\theta)$ is motivated by the observation that if
$\expect{}{\yn | \z_m, \theta} = \expect{}{\yn | \z_m}$ at our particular
$\theta$, then $\g(\theta) \approx \expect{\target(\yn)}{\yn}$ by the law of
large numbers, since $\z_m$ are drawn marginally from the target distribution of
regressors.  The Bayesian estimator integrates out our uncertainty in the
$\theta$, taking $\expect{\post}{\g(\theta)}$ as our posterior estimate of the
unknown $\expect{\target(\yn)}{\yn}$.  See
\citet{park:2004:mrp,lopezmartin:2022:mrptutorial} for more details.

To assess our ability to estimate $\gcovtrue$ in MrP models, we replicate a
semi--simulated analysis from Chapter 1 of the MrP tutorial textbook,
\citet{lopezmartin:2022:mrptutorial}.  Although this MrP analysis is based on
real data, the survey population consists of a sample taken uniformly at random
from a \emph{superpopulation} for which complete data --- both regressors and
responses --- is available.  The MrP estimator $\expect{\post}{\g(\theta)}$,
which is formed using response variables only from the sampled survey population,
can then be compared to the known ground truth: namely, the average response in
the superpopulation. \citet{lopezmartin:2022:mrptutorial} use the
superpopulation to assess the accuracy of MrP methods.  In our case, we use
the superpopulation to estimate the ground truth $\gcovtrue$ by measuring the
sampling variability of $\expect{\post}{\g(\theta)}$ under repeated independent
draws of the survey population.

The superpopulation dataset provided by \citet{lopezmartin:2022:mrptutorial}
consists of $\mrpNTotalobs$ pairs of survey responses and regressors, from which
we sample $\mrpNobs$ uniformly at random without replacement to form simulated
survey populations.  The response $\yn$ is an indicator of whether a respondent
supports coverage of abortions in insurance plans, taken from the 2018
Cooperative Congressional Election Study \citep{schaffner:2018:cces}.  The
provided regressors, $\zn$, consist of discretized demographic information
measuring age, gender, ethnicity, education, and state.  We use the model given
by Chapter 1 of \citet{lopezmartin:2022:mrptutorial}, which consisted of a
hierarchical logistic regression with unknown hyperparameter variances.  We let
our parameter $\theta$ include both logistic regression coefficients and
hyperprior variances; in total, the dimension of the parameter space was
$\thetadim = \mrpThetaDim$.

The details of our sampling and estimation procedure were as follows.  A draw of
the survey data $\xvec$ is produced by selecting $\mrpNobs$ regressor / response
pairs $\x_n = (\z_n, \y_n)$ uniformly at random without replacement from the
$\mrpNTotalobs$ observations in the superpopulation.  For each such $\xvec$, we
can compute $\post$ and so $\expect{\post}{\g(\theta)}$, using the
$\mrpNTotalobs - \mrpNobs$ remaining superpopulation observations as the target
population in $\g(\theta)$.\footnote{Technically there will be additional
variability due to the random selection of the target regressors in
\eqref{mrpg}, but we are justified in neglecting this variability since the
superpopulation is so large.} By simulating multiple draws of $\xvec$ and
computing $\expect{\post}{\g(\theta)}$ for each, we can estimate the ground
truth $\gcovtrue$.   In this way, we simulated $\mrpNSamples$ different survey
datasets $\xvec$.  For each dataset, we drew $M=\mrpNMCMC$ MCMC samples from
each posterior using the \texttt{R} package \texttt{brms}
\citep{burkner:2018:brms}. We then used these samples to compute an MCMC
estimate of $\expect{\post}{\g(\theta)}$ on each, the variance of which across
distinct simulated values of $\xvec$ provides a ground--truth estimate of
$\gcovtrue$.  On each set of MCMC draws, we computed $\gcovijhat$ and
$\var{\post}{\g(\theta)}$.  We randomly chose one of the samples as a basis for
$\mrpNBoot$ bootstrap replicates to compute a bootstrap variance estimate
$\gcovboothat$.

\MRPGraph{}

As shown in \figref{mrp_graph}, the IJ covariance and bootstrap covariance are
nearly identical.  In this case, the Bayesian posterior variance also coincides
with the frequentist variance.  The total computation time for all bootstrap
samples was $\mrpBootMinutes$ minutes, and the computation for the IJ covariance
on the same set of MCMC draws used for the bootstraps was $\mrpIJSec$
seconds.\footnote{ The computation time for the IJ covariance was dominated by a
single call to \texttt{brms::posterior\_linpred}, which took $\mrpLinpredSec$ of
the $\mrpIJSec$ seconds. The function \texttt{brms::posterior\_linpred} is only
computing the log odds for each observation and each MCMC draw, but is doing so
with a sequence of nested loops rather than vectorized operations, due to the
idiosyncrasies of \texttt{rstan} data objects.  In a different MCMC environment,
this computation could presumably be done much more efficiently.}

As we show in \appref{experiments_mrp_appendix}, on this problem the MAP
estimate provides a poor estimate of the posterior expectation, likely due to
the relatively high dimensionality of $\theta$ and the nonlinearity of the
model, particularly through estimation of the random effect variances.
Nevertheless, since $N / \thetadim \approx \mrpNperD$, it is plausible that the
posterior is concentrated enough in the nuisance parameters to avoid the
challenges for IJ in the high--dimensional regime we discussed in
\secref{data_weights}. In \appref{experiments_mrp_appendix}, we also compare
with marginal maximum likelihood estimates provided by the \texttt{R} function
\texttt{lmer} \citep{bates:2015:lme4}, noting that \texttt{lmer} does not
naturally provide estimates of $\gcovtrue$ for a complicated quantity depending
jointly on the random and fixed effects.



